\documentclass[11pt]{article}

\usepackage[margin=1in]{geometry}
\usepackage{amsmath,amssymb,amsthm}
\usepackage{array}
\usepackage{booktabs}
\usepackage{microtype}
\usepackage[hidelinks]{hyperref}

\newtheorem{theorem}{Theorem}
\newtheorem{proposition}[theorem]{Proposition}
\newtheorem{lemma}[theorem]{Lemma}
\newtheorem{corollary}[theorem]{Corollary}
\theoremstyle{definition}
\newtheorem{definition}[theorem]{Definition}
\theoremstyle{remark}
\newtheorem{remark}[theorem]{Remark}

\newcommand{\cube}{\{0,1\}}
\newcommand{\F}{\mathbb F}
\newcommand{\dist}{d_{\mathrm H}}
\newcommand{\K}{K_2(11,3)}

\title{Certified Branch Reductions for the Binary Covering Number
  $K_2(11,3)$}
\author{Ruturaj R Raval\\
  Independent Researcher\\
  \href{https://orcid.org/0000-0003-4930-8981}
  {ORCID: 0000-0003-4930-8981}}
\date{September 2026}

\begin{document}
\maketitle

\begin{abstract}
The binary covering number $K_2(11,3)$ is recorded in the interval
$15\leq K_2(11,3)\leq16$.  We give a reproducible, certificate-oriented
reduction of the unresolved size-15 search.  Exact shell and
parity-refined Delsarte certificates show that every hypothetical
15-word cover has minimum distance at most 5, total pair-ball overlap at
least 1712, and total triple-ball overlap at least 280.  Translation and
coordinate symmetry then yield a complete 150-case normalized
partition.  Exact orbit-profile certificates, integer-profile
enumeration, checked DRAT traces, and maximum-degree normalization close
112 normalized branches and leave 38 explicit residual branches.
Those residual branches contain 2548 third-word orbit children before
static exclusions and 2163 live children afterward.  In four selected
unresolved children, an exhaustive 350-branch fourth-word split has 184
reverse-unit-propagation closures and 140 additional solver-generated
DRAT closures.  Thus 324 selected branches are certified closed and 26
remain.  No complete selected child or normalized parent is closed by
this final layer.  The exact value and the known interval are unchanged.
\end{abstract}

\section{Introduction}

A binary code $C\subseteq\F_2^n$ has covering radius at most $R$ when
every word in $\F_2^n$ is within Hamming distance $R$ of some codeword.
The covering number $K_2(n,R)$ is the minimum possible size of such a
code.  Covering codes are a classical subject in coding theory, with
connections to data compression, test design, combinatorial coverings,
and domination in Hamming graphs \cite{cohen-book,graham-sloane}.

The binary covering-code table maintained by Keri records
\[
  15\leq K_2(11,3)\leq16.
\]
Its update log dates the lower-bound entry to 2006, while the table
attributes the upper-bound line to the Graham--Sloane construction
record \cite{keri-table,graham-sloane}.  A targeted literature audit in
September 2026 found no later primary source determining this exact
cell.  This dated audit is not a proof of absence.  Recent work on
semidefinite lower bounds also illustrates that a modern,
proof-carrying treatment of individual table cells remains useful
\cite{gijswijt-polak}.

There are two definitive ways to close the gap.  A verified 15-word
cover would prove $K_2(11,3)=15$.  A checked exclusion of every
15-word candidate, combined with a verified 16-word cover, would prove
$K_2(11,3)=16$.  The work reported here does neither.  Its contribution
is a finite and reproducible reduction of the second route.

\begin{table}[ht]
\centering
\begin{tabular}{@{}ll@{}}
\toprule
Item & Status \\
\midrule
Known interval & $15\leq K_2(11,3)\leq16$ \\
Verified 16-word radius-3 cover & Yes \\
Verified 15-word radius-3 cover & No \\
Complete normalized parent partition & 150 branches \\
Certified normalized branch closures & 112 \\
Residual normalized branches & 38 \\
Selected fourth-word branches & 350 \\
Certified selected fourth-word closures & 324 \\
Exact value determined & No \\
\bottomrule
\end{tabular}
\caption{The certified frontier.  The fourth-word layer is local to
four selected children and does not reduce the global count of 38
residual normalized branches.}
\label{tab:frontier}
\end{table}

The reduction combines exact analytic inequalities, Hamming-cube
automorphisms, orbit averaging, exact integer enumeration, and
proof-producing satisfiability solving.  Every retained SAT-based
exclusion is tied to a reconstructed formula and a separately
checked proof trace.  The complete proof artifacts are archived with
the version-specific DOI
\href{https://doi.org/10.5281/zenodo.22302261}
{10.5281/zenodo.22302261}.  The archived file and its SHA-256 digest
are, respectively,
\begin{center}
\begin{minipage}{0.94\linewidth}
\small\ttfamily
ruturajr-raval/binary-covering-code-11-3-v0.2.0.zip\\
750003eba2e9f9baf5fee9ed93c679b3661daf6d8c68ca40eeb681202b5e72ff
\end{minipage}
\end{center}
After the source archive is authenticated against the SHA-256 published
with the release, its compact replay checks the principal ledgers, the
root and fourth-word proof manifests, and the Zenodo archive binding.
It also regenerates the analytic certificates and verifies the known
16-word construction.

\section{Definitions and the verified upper endpoint}

\begin{definition}
For $x,y\in\F_2^n$, let $\dist(x,y)$ be their Hamming distance.  The
radius-$R$ ball centered at $x$ is
\[
  B_R(x)=\{y\in\F_2^n:\dist(x,y)\leq R\}.
\]
A code $C\subseteq\F_2^n$ is an $(n,R)$ covering code when
\[
  \bigcup_{c\in C} B_R(c)=\F_2^n.
\]
\end{definition}

For $(n,R)=(11,3)$, the ambient space has 2048 words and each ball has
\[
  |B_3(0)|=\sum_{i=0}^{3}\binom{11}{i}=232.
\]
The repository contains the following 16-word linear code:
\[
\begin{array}{llll}
00000000000 & 00010001001 & 00101010101 & 00111011100\\
01001101011 & 01011100010 & 01100111110 & 01110110111\\
10001001000 & 10011000001 & 10100011101 & 10110010100\\
11000100011 & 11010101010 & 11101110110 & 11111111111.
\end{array}
\]
Direct enumeration and a syndrome-space computation both give covering
radius exactly 3.  This verifies the known upper bound
$K_2(11,3)\leq16$; no originality is claimed for the construction.

\section{Exact structural constraints at size 15}

Assume throughout this section that $C\subseteq\F_2^{11}$ is a
radius-3 cover with $|C|=15$.  Let $p_d$ be the number of unordered
codeword pairs at Hamming distance $d$.

\subsection{Shell and Delsarte rows}

There are 105 unordered pairs, so
\[
  \sum_{d=1}^{11}p_d=105.
\]
For a fixed codeword and a target shell at distance $j$, define
\[
  H(d,j)=
  \sum_{\substack{0\leq a\leq d\\
                  0\leq j-a\leq 11-d\\
                  d+j-2a\leq3}}
    \binom da\binom{11-d}{j-a}.
\]
Summing the coverage requirement over codewords and over every target
in shell $j$ gives, for $4\leq j\leq11$,
\[
  2\sum_{d=1}^{11}H(d,j)p_d
  \geq15\binom{11}{j}.
\]

Let $K_k(d)$ be the binary Krawtchouk polynomial.  Fourier expansion of
the distance distribution gives
\[
  15\binom{11}{k}
  +2\sum_{d=1}^{11}K_k(d)p_d
  =
  \sum_{\operatorname{wt}(a)=k}
  \left(\sum_{c\in C}(-1)^{a\cdot c}\right)^2.
\]
Every inner sum is odd because $|C|=15$.  Its square is therefore at
least 1, yielding the parity-refined inequality
\[
  2\sum_{d=1}^{11}K_k(d)p_d
  \geq-14\binom{11}{k}.
\]

\begin{theorem}[Distance-distribution constraints]
\label{thm:distance}
Every 15-word radius-3 cover in $\F_2^{11}$ satisfies
\[
  \sum_{d=1}^{6}p_d\geq28,
  \qquad
  \sum_{d=1}^{5}p_d\geq11.
\]
In particular, its minimum distance is at most 5.
\end{theorem}

\begin{proof}
The first inequality is the exact nonnegative combination
\[
  \tfrac13 P+\tfrac1{24}D_1+\tfrac1{24}D_{11},
\]
where $P$ is the pair-count equality and $D_k$ is the degree-$k$
parity-refined Delsarte row.  Its coefficient vector, ordered by
distances 1 through 11, is
\[
  (1,1,\tfrac23,\tfrac23,\tfrac13,\tfrac13,
    0,0,-\tfrac13,-\tfrac13,-\tfrac23).
\]
This is coefficientwise at most the indicator of distances 1 through 6,
and the resulting rational lower bound is 28.

For the second inequality, combine
\[
  \tfrac{37}{615}P+\tfrac1{410}S_4
  +\tfrac{11}{82}S_{11}
  +\tfrac{33}{820}D_1+\tfrac{13}{2460}D_9,
\]
where $S_j$ denotes the shell-$j$ row.  Coefficientwise comparison with
$p_1+\cdots+p_5$ gives the lower bound $12467/1230$.  The left side is
an integer, so it is at least 11.  Both combinations and every
coefficient are recomputed using exact rational arithmetic by the
ancillary verifier.
\end{proof}

\subsection{Pair and triple overlap}

Let
\[
  O=\sum_{\{c,c'\}\subseteq C}|B_3(c)\cap B_3(c')|
\]
be the total pair-ball overlap.  The intersection sizes for pair
distances $1,\ldots,11$ are
\[
  112,112,56,56,20,20,0,0,0,0,0.
\]
Normalize the shell-10, shell-11, degree-1 Delsarte, and degree-11
Delsarte rows by dividing their even coefficients by 2.  Explicitly,
\[
\begin{array}{rcl}
J_{10}&=&(0,0,0,0,0,0,4,3,11,11,11)\cdot p\geq83,\\
J_{11}&=&(0,0,0,0,0,0,0,1,1,1,1)\cdot p\geq8,\\
D_1&=&(9,7,5,3,1,-1,-3,-5,-7,-9,-11)\cdot p\geq-77,\\
D_{11}&=&(-1,1,-1,1,-1,1,-1,1,-1,1,-1)\cdot p\geq-7.
\end{array}
\]
Here $p=(p_1,\ldots,p_{11})$.  Exact coefficient comparison gives
\[
\begin{split}
  O={}&20P+4J_{10}+4J_{11}+9D_1+9D_{11}\\
     &+20(p_1+p_2)+4(p_9+p_{10})+40p_{11}.
\end{split}
\]
The row bounds give a base value 1708.  Every coefficient of
$J_{10}+J_{11}$ is divisible by 4, while $83+8\equiv3\pmod4$.
The shell slacks $u=J_{10}-83$ and $v=J_{11}-8$ therefore satisfy
$u+v\equiv1\pmod4$, hence $u+v\geq1$.

\begin{theorem}[Overlap constraints]
\label{thm:overlap}
Every 15-word radius-3 cover satisfies
\[
  O\geq1712+20(p_1+p_2)+4(p_9+p_{10})+40p_{11}.
\]
In particular, $O\geq1712$.  If
\[
  T=\sum_{\{c_1,c_2,c_3\}\subseteq C}
  |B_3(c_1)\cap B_3(c_2)\cap B_3(c_3)|,
\]
then $T\geq280$.
\end{theorem}

\begin{proof}
The displayed exact identity and the shell-slack congruence prove the
pair-overlap inequality.  For the triple term, let $m_x$ be the number
of codeword balls containing $x$.  Since $C$ is a cover, $m_x\geq1$.
Pointwise,
\[
  \binom{m_x}{3}\geq\binom{m_x}{2}-(m_x-1).
\]
The total incidence excess is
\[
  15\cdot232-2048=1432.
\]
After summing over $x\in\F_2^{11}$, we obtain
$T\geq O-1432\geq280$.
\end{proof}

\section{The size-15 formula}

Introduce one Boolean variable $x_v$ for each $v\in\F_2^{11}$.
Coverage is expressed by the 2048 clauses
\[
  \bigvee_{\dist(u,v)\leq3}x_v,
  \qquad u\in\F_2^{11}.
\]
An exact or at-most-15 cardinality constraint restricts the selected
codewords.  For a nonempty covering code, an at-most-15 model can be
padded to exactly 15 distinct codewords without losing coverage, so the
two formulations are equivalent for existence.

Two separately audited encodings were retained.  A sequential
exact-cardinality formula has 34816 variables and 131029 clauses.  A
compact at-most-15 totalizer formula has 9464 variables and 26804
clauses.  Projection tests on smaller instances verify that auxiliary
variables preserve the intended codeword semantics.

The Hamming cube automorphism group contains all translations and
coordinate permutations.  The reduction uses only transformations that
are explicitly shown to preserve coverage and the formula assumptions.
No solver-discovered symmetry-breaking condition enters the proof
ledger without a separate completeness argument.

\section{A complete 150-case normalized partition}

Since $15\cdot232>2048$, the 15 radius-3 balls cannot be pairwise
disjoint.  Thus some pair of codewords has distance at most 6.  Choose
a globally closest pair, let its distance be $w$, and translate one
endpoint to zero.  The other endpoint is then a minimum-weight nonzero
codeword of the translated code.  A coordinate permutation maps it to
\[
  A=2^w-1,\qquad 1\leq w\leq6.
\]
The stabilizer of $A$ is $S_w\times S_{11-w}$.  Its orbit on a second
word $B$ is determined by
\[
  \bigl(\operatorname{wt}(B),
        |\operatorname{supp}(A)\cap\operatorname{supp}(B)|\bigr).
\]
Choose the lexicographically first occupied orbit and map one selected
word in it to its canonical representative.  Across the six possible
minimum weights, this produces 150 canonical parent branches.

\begin{theorem}[Certified normalized frontier]
\label{thm:parents}
If a 15-word radius-3 cover exists, a Hamming-cube automorphism and the
documented closest-pair and maximum-degree normalizations map it into
one of 38 explicitly retained normalized branches.
\end{theorem}

\begin{proof}
The complete two-word orbit manifest contains 150 branches.  The
stage-1 ledger classifies them in the disjoint precedence order shown
in Table~\ref{tab:parents}.  Exact Farkas multipliers close 80 branches.
Exhaustive integer orbit-profile certificates close two more.  One
standalone formula and 14 closest-pair formulas have checked DRAT
proofs.  Theorem~\ref{thm:distance} closes the four remaining
weight-6 branches in the broader first-stage ledger.

The 49 stage-1 residual branches are then normalized as follows.  For a
code of minimum distance $d$, form the graph on its codewords in which
two vertices are adjacent exactly when their distance is $d$.  Choose a
maximum-degree vertex, translate it to zero, choose one of its neighbors,
and map that neighbor to $A=2^d-1$.  This is compatible with the
two-word partition because translations and coordinate permutations
preserve distances and coverage.

Let $B$ be the selected word in the first remaining descriptor.  If
$\operatorname{wt}(B)=d$, then zero has at least two minimum-distance
neighbors, and no additional conclusion is drawn.  If
$\operatorname{wt}(B)>d$, the descriptor ordering implies that $A$ is
the only selected word of weight $d$, so zero has degree one.  If
$\dist(A,B)=d$, then $A$ has at least the two neighbors zero and $B$,
contradicting the maximum-degree choice.  This excludes five branches.
If instead $\dist(A,B)>d$, the maximum degree is one and the entire
minimum-distance graph is a matching.  This proves the completeness of
the classification into 5 contradiction, 34 matching, and 10
multiple-neighbor cases.

Six complete third-word formulas have checked DRAT proofs, four within
the separately proved matching scope and two without that condition.
The machine ledger checks disjointness, all 150 identities, all input
hashes, and the final count $150-112=38$.
\end{proof}

\begin{table}[ht]
\centering
\begin{tabular}{@{}lr@{}}
\toprule
Closure reason & Branches \\
\midrule
Exact rational orbit-profile infeasibility & 80 \\
Exact integer orbit-profile infeasibility & 2 \\
Standalone checked DRAT & 1 \\
Minimum distance at most 5 & 4 \\
Closest-pair checked DRAT & 14 \\
Maximum-degree contradiction & 5 \\
Checked third-word DRAT & 6 \\
\midrule
Certified normalized closures & 112 \\
Residual normalized branches & 38 \\
\bottomrule
\end{tabular}
\caption{The complete normalized parent ledger.  The advanced rows are
normalization statements and must not be reinterpreted as 11 additional
unrestricted parent-CNF exclusions.}
\label{tab:parents}
\end{table}

\begin{remark}
The exact rational orbit analysis leaves 70 cases feasible over the
reals.  Integer orbit counts exclude two of these, leaving 68 feasible
aggregate profiles.  Such profiles are necessary averaged data, not
vertex-level covering codes.
\end{remark}

\section{Third-word and fourth-word orbit frontiers}

Fix a canonical parent with selected words $0,A,B$.  The ordered pair
$(A,B)$ partitions the 11 coordinates into four cells:
\[
  A\cap B,\quad A\setminus B,\quad B\setminus A,\quad
  [11]\setminus(A\cup B).
\]
The pointwise stabilizer permutes coordinates within these cells.
Therefore the orbit of a third word $D$ is determined by its four
intersection counts.  Selecting the first occupied orbit gives an
exhaustive child split.

The 49 stage-1 residual parents have 3238 such child orbits.  Six
third-word DRAT closures remove 423 children, and five
maximum-degree-closed parents account for 267 more.  The 38 active
parents contain 2548 raw children.  Existing minimum-distance clauses
exclude 312, and the certified matching condition excludes 73
additional children in its valid scope.  Exactly 2163 live children
remain.

The next retained layer studies four selected unresolved children below the parent
\texttt{w4-weight5-intersection0}.  Once three nonzero words are fixed,
the coordinates split into eight membership cells.  A fourth word is
classified by its eight intersection counts.  In each selected child,
the first occupied fourth-word orbit gives an exhaustive split because
the fixed words leave more than 1200 ambient vertices uncovered.

\begin{theorem}[Selected fourth-word reduction]
\label{thm:fourth}
The four selected third-word children in Table~\ref{tab:fourth} split
into 350 exhaustive fourth-word branches.  Checked reverse unit
propagation certificates close 184 branches, and checked
solver-generated DRAT certificates close 140 more.  Exactly 26 branches
remain unresolved.  Every selected child retains at least one residual
branch.
\end{theorem}

\begin{proof}
The separate orbit auditor reconstructs the eight coordinate cells,
all candidate words, static exclusions, canonical representatives, and
branch digests.  It obtains 350 branches.

For 184 formulas, the branch units force a contradiction by reverse
unit propagation.  Each resulting trace is checked against the
reconstructed formula.  The remaining 166 formulas formed the
exploratory solver frontier.  A retained proof-producing rerun generated
DRAT traces for 140 of them, while 26 remained unresolved.  Each
retained trace is checked by \texttt{drat-trim} at pinned revision
\texttt{2e3b2dc0ecf938addbd779d42877b6ed69d9a985}.  The proof indexes
authenticate every formula, trace, checker result, and production
record.  Clean-checkout self-attestation replay passed for both layers.
The per-child residual counts in Table~\ref{tab:fourth} sum to 26.
\end{proof}

\begin{table}[ht]
\centering
\begin{tabular}{@{}lrrrr@{}}
\toprule
Child & RUP & DRAT & Residual & Total \\
\midrule
\texttt{orbit-005} & 50 & 29 & 6 & 85 \\
\texttt{orbit-007} & 53 & 15 & 8 & 76 \\
\texttt{orbit-014} & 41 & 28 & 4 & 73 \\
\texttt{orbit-015} & 40 & 68 & 8 & 116 \\
\midrule
Total & 184 & 140 & 26 & 350 \\
\bottomrule
\end{tabular}
\caption{Certified branch-level results inside four selected unresolved
third-word children.}
\label{tab:fourth}
\end{table}

\begin{corollary}
The selected fourth-word layer closes no complete third-word child and
no normalized parent.  Consequently, it does not reduce the global
count of 38 residual normalized branches.
\end{corollary}

\section{Certificate architecture and reproduction}

The proof records separate mathematical scope from computational
status.  A timeout, an unlogged UNSAT response, or a heuristic search
failure is never entered as an exclusion.  A retained SAT-based closure
contains:
\begin{enumerate}
  \item an exact formula identity and semantic metadata;
  \item a deterministic formula reconstruction path;
  \item a SHA-256 digest of the formula;
  \item a compressed RUP or DRAT trace;
  \item a checker record tied to a pinned \texttt{drat-trim} revision;
  \item an exact-membership manifest for the proof bundle; and
  \item a clean-checkout replay record.
\end{enumerate}
The DRAT format permits compact separately checked certificates for
SAT-solver refutations \cite{heule-drat}.

The arXiv source archive intentionally omits the large proof payloads.
It contains the exact ledgers, proof indexes, replay records, the
16-word code, and a deterministic manifest.  Running
\begin{verbatim}
python3 replay.py
\end{verbatim}
from the extracted source:
\begin{enumerate}
  \item checks every ancillary file against the internal manifest;
  \item verifies the 16-word covering radius by direct enumeration;
  \item checks the recorded 150-to-38 and 350-to-26 ledger counts,
        statuses, and source digests;
  \item checks the retained RUP and solver-DRAT proof-index identities;
  \item checks the retained clean-checkout replay statuses; and
  \item regenerates the distance and overlap certificates
        byte-for-byte using standard-library Python.
\end{enumerate}
No network access or third-party package is required for this compact
replay.  Full replay of all 324 retained proof traces requires the
immutable v0.2.0 artifact deposit at
\href{https://doi.org/10.5281/zenodo.22302261}
{10.5281/zenodo.22302261} and the pinned verification environment.
The all-versions concept DOI is
\href{https://doi.org/10.5281/zenodo.22260709}
{10.5281/zenodo.22260709}.

\section{Limitations and next steps}

The known interval remains
\[
  15\leq K_2(11,3)\leq16.
\]
The present work does not prove a new lower bound and does not produce
a 15-word cover.  Thirty-eight normalized parent branches remain
unresolved.  The 2163 live third-word children have not been
exhaustively decided.  The selected fourth-word bundle leaves 26
branches and closes no complete selected child.

The immediate proof route is an audited fifth-word orbit split of those
26 branches.  Closing all four selected children would still leave
other live children below the 38 residual parents, so subsequent work
must expand the certified cover rather than extrapolate from local
success.  In parallel, nonlinear construction search remains valuable:
a single independently verified 15-word cover would settle the problem
in the opposite direction.

The methods are reusable beyond this table cell.  Exact orbit averaging,
integer profile certificates, stabilizer-aware SAT formulas, and
proof-logged branch exclusions apply to finite covering and domination
problems with large automorphism groups.

\section{Conclusion}

We have converted a broad size-15 search into an explicit certified
frontier.  Every hypothetical cover can be normalized into one of 38
retained branches.  Within four selected unresolved children, 324 of 350
fourth-word branches have checked certificates.  These are substantial
branch-level reductions, but they are not a solution of $K_2(11,3)$.
The remaining cases are stated explicitly so that future computation
can extend a verified proof cover rather than repeat an unauthenticated
search.

\begin{thebibliography}{9}

\bibitem{cohen-book}
G.~Cohen, I.~Honkala, S.~Litsyn, and A.~Lobstein,
\emph{Covering Codes},
North-Holland, 1997.

\bibitem{gijswijt-polak}
D.~Gijswijt and S.~Polak,
``Semidefinite lower bounds for covering codes,''
arXiv:2504.01932, version 2, 2026.

\bibitem{graham-sloane}
R.~L. Graham and N.~J.~A. Sloane,
``On the covering radius of codes,''
\emph{IEEE Transactions on Information Theory},
31(3):385--401, 1985.
doi:10.1109/TIT.1985.1057039.

\bibitem{heule-drat}
M.~J.~H. Heule,
``The DRAT format and DRAT-trim checker,''
arXiv:1610.06229, 2016.

\bibitem{keri-table}
G.~Keri,
``Tables of covering codes,''
\url{https://old.sztaki.hu/~keri/codes/}.
Accessed September 2026.

\end{thebibliography}

\end{document}
